% continuum-measure.tex — Chapter on the projective-limit construction
% of the continuum SU(2) Wilson measure mu_infinity.
%
% Linked code / Lean:
%   - computations/continuum_limit.py (numerical UV-trajectory data)
%   - lean/YMGap/Continuum/ContinuumMeasure.lean (type signature only;
%     analytical existence lives at paper level)
%   - Milestone 7.2 of work_plan.md.

\chapter{The Continuum
  \texorpdfstring{$\SUtwo$}{SU(2)} Wilson Measure
  \texorpdfstring{$\muinf$}{mu\_infinity}:
  Projective-Limit Construction}
\label{ch:continuum-measure}

\noindent\emph{Daniell--Kolmogorov extension of the UV-direction
RG-flow lattice measures to a continuum cylindrical probability
measure $\muinf$ on $\Ominf$.}

We construct the continuum $\SUtwo$ Wilson measure $\muinf$ as the
projective limit of the finite-volume lattice measures
$\{\muk\}_{k \ge 0}$ generated by the UV-direction block-spin flow
of Milestones~6.3 and~7.1.  The construction proceeds in four steps:
(i)~assemble the lattice measures into a projective system indexed
by the UV-direction RG flow with transition maps $\Tbstar$;
(ii)~verify the Kolmogorov consistency condition $\Tbstar \muk =
\mu_{k+1}$; (iii)~establish tightness via the compactness of the
per-link target $\SUtwo^{\Ek}$; (iv)~invoke the Daniell--Kolmogorov
extension theorem to produce a unique cylindrical probability measure
$\muinf$ on the projective-limit space $\Ominf$, into which the
Phase-7.1 limit $\sighinf > 0$ injects as the area-law-decay rate of
Wilson-loop expectations under $\muinf$.

This chapter is the paper-level deliverable of Milestone~7.2.  Per
the work-plan, the analytical existence proof lives at the paper
level (the projective-limit construction for compact-group lattice
measures is standard but heavier than the $\Qp$~/~constructive-reals
stack currently in Lean); the Lean module
\codepath{lean/YMGap/Continuum/ContinuumMeasure.lean} provides only
the type signature, with the existence claim referenced back here.

\section{Introduction and standing setup}
\label{sec:continuum-intro}

\subsection{Standing data from upstream milestones}
\label{subsec:upstream}

We work over the four-dimensional Euclidean lattice $\mathbb{Z}^4$ at
scale~$k$ with lattice spacing $a_k$, link set $\Ek$, and per-link
target $SU(2)$.  The basic configuration space at scale~$k$ is
\begin{equation*}
  \Omk \;=\; SU(2)^{\Ek},
\end{equation*}
equipped with the product Haar measure as reference; this is a compact
Polish space (a~finite product of compact Polish factors).  The
finite-volume Wilson measure $\muk$ on $\Omk$ is the Gibbs measure of
the Wilson action at coupling $g_k$:
\begin{equation*}
  \muk(\mathrm{d}U) \;=\; Z_k^{-1} \, \exp\bigl(- S_k^{\mathrm{Wilson}}(U)\bigr) \, \mathrm{d}U,
\end{equation*}
where the action sums over plaquettes of the scale-$k$ lattice with
the standard heat-kernel weight.  Per
\texttt{computations/uv\_trajectory.py} and
\texttt{formal/block-spin-morphism.tex} we move \emph{UV-ward} by
refining ($a_{k+1} = a_k / b$, $b > 1$); the coupling evolves as
$g_{k+1}$ via the one-loop SU(2) $\beta$-function
\begin{equation*}
  \frac{1}{g_{k+1}^2}  \;=\;  \frac{1}{g_k^2} \;+\; \frac{b_0}{16\pi^2} \, \log b,
  \qquad
  b_0 \;=\; 22/3,
\end{equation*}
and remains strictly inside the small-field convergent regime.

The block-spin map $\Tb \colon \Omkp \to \Omk$ of
\texttt{formal/block-spin-morphism.tex} Definition~3.1, read in the
UV direction, is a measurable \emph{coarsening} that sends fine-lattice
configurations to coarse-lattice configurations by averaging or
projection.  The pushforward $\Tbstar = (\Tb)_*$ takes a probability
measure on $\Omkp$ to a probability measure on $\Omk$.

\begin{convention}[UV indexing]
\label{conv:uv}
Throughout this document, the index $k$ \emph{increases} along the
UV-direction RG flow.  Hence $a_k$ decreases in $k$,
$\dim \Omk = |\Ek|$ increases in $k$, the coupling $g_k$ decreases in $k$
toward zero, and the lattice $\Omk$ becomes finer.  The projective
limit is over the cofiltered system
\begin{equation*}
  \cdots \xleftarrow{\Tb} \Omega_2 \xleftarrow{\Tb} \Omega_1 \xleftarrow{\Tb} \Omega_0,
\end{equation*}
where every arrow is a coarsening and the limit object $\Ominf$ is the
\emph{finest} continuum-like configuration space.
\end{convention}

\subsection{Scope and what this document does not do}
\label{subsec:scope}

This document constructs $\muinf$ as a cylindrical probability measure
on $\Ominf$ via the Daniell--Kolmogorov projective-limit theorem.  It
does \emph{not}:

\begin{itemize}
  \item Identify $\muinf$ with any continuum quantum field-theoretic
    measure beyond the projective-limit construction.  The
    identification with the OS-axiom-compliant continuum Wilson
    measure is the deliverable of Phase~8.1.
  \item Establish the Osterwalder--Schrader axioms for $\muinf$
    (Phase~8.1).
  \item Lift the Lean formalization beyond a type signature.  The
    analytical content (consistency, tightness, the extension
    theorem) lives here.  Per the work-plan Milestone~7.2 open
    question, this is the deliberate decision point at which $Q'$
    no longer suffices; the project has chosen the \textbf{paper-only}
    branch and deferred any \texttt{Constructive/Reals.lean} build-out
    to Milestone~7.3 (which is downstream of the mass-gap statement
    and is not on the critical path for stating the result).
\end{itemize}

\section{The projective system $\{\muk, \Tbstar\}$}
\label{sec:projective-system}

\subsection{The cofiltered index category}
\label{subsec:cofiltered}

\begin{definition}[UV index category]
\label{def:uv-index}
Let $\mathbf{UV}$ be the small category whose objects are the
nonnegative integers $\mathbb{N}$ and whose unique morphism
$k+1 \to k$ is the block-spin coarsening $\Tb$.  $\mathbf{UV}$ is
the linear poset $(\mathbb{N}, \ge)$ viewed as a category; it is
\emph{cofiltered} (every pair of objects has a common predecessor:
$\max(j, k) \to j$ and $\max(j, k) \to k$).
\end{definition}

The opposite category $\mathbf{UV}^{\mathrm{op}}$ is filtered.
Sequences indexed by~$\mathbf{UV}$ form an inverse system; we will
take their inverse limit (= projective limit) in the categories of
compact Hausdorff spaces and of probability measures.

\subsection{The configuration projective system}
\label{subsec:config-system}

\begin{definition}[Configuration projective system]
\label{def:config-system}
The \emph{configuration projective system} is the functor
\begin{equation*}
  \Omega_\bullet \;\colon\; \mathbf{UV} \;\longrightarrow\; \mathbf{CHaus}
\end{equation*}
that sends the object $k$ to $\Omk = SU(2)^{\Ek}$ and sends the
morphism $k+1 \to k$ to $\Tb \colon \Omkp \to \Omk$.  Functoriality
holds because consecutive coarsenings compose, $\Tb \circ \Tb$ is
itself a coarsening at block-factor $b^2$ (consistent with the
$\Tb^n$ notation used in \texttt{formal/block-spin-morphism.tex}~\S4
for iterated steps).
\end{definition}

\begin{definition}[The projective-limit configuration space]
\label{def:proj-limit-config}
The \emph{continuum configuration space} is the projective limit
\begin{equation*}
  \Ominf \;:=\; \varprojlim_{k \in \mathbf{UV}^{\mathrm{op}}} \, \Omk,
\end{equation*}
a compact Hausdorff space concretely realized as the closed subset
\begin{equation*}
  \Ominf \;=\; \Bigl\{ \, (U_k)_{k \ge 0} \in \prod_k \Omk
                     \;:\;
                     \Tb(U_{k+1}) = U_k \;\; \forall k \, \Bigr\}
              \;\subseteq\; \prod_k \Omk.
\end{equation*}
By Tychonoff's theorem, $\prod_k \Omk$ is compact; the inverse-system
constraint is closed, so $\Ominf$ is itself compact Hausdorff (in fact
Polish, being a closed subset of a countable product of compact metric
spaces).
\end{definition}

\begin{proposition}[Projection maps]
\label{prop:projections}
For each $k$, the natural projection
$\pi_k \colon \Ominf \to \Omk$ given by $(U_j)_{j \ge 0} \mapsto U_k$
is continuous and surjective.  Moreover $\Tb \circ \pi_{k+1} = \pi_k$
by definition of $\Ominf$.
\end{proposition}

\subsection{The measure projective system}
\label{subsec:measure-system}

\begin{definition}[Measure projective system]
\label{def:measure-system}
The \emph{measure projective system} is the functor
\begin{equation*}
  \mu_\bullet \;\colon\; \mathbf{UV} \;\longrightarrow\; \mathbf{Prob}
\end{equation*}
sending $k \mapsto \muk$ (the Wilson measure at scale~$k$) and
the morphism $k+1 \to k$ to $\Tbstar \colon \mathbf{Prob}(\Omkp) \to
\mathbf{Prob}(\Omk)$.  Functoriality is automatic from functoriality
of $\Omega_\bullet$ and the standard pushforward
$\Tbstar \mu(E) := \mu(\Tb^{-1}(E))$.
\end{definition}

We do not (yet) claim functoriality means the system is
\emph{consistent} in the Kolmogorov sense; that is the content of
\S\ref{sec:consistency}.

\section{The Kolmogorov consistency condition}
\label{sec:consistency}

\begin{definition}[Consistency]
\label{def:consistency}
A family $\{\nu_k\}_{k \ge 0}$ of probability measures with
$\nu_k \in \mathbf{Prob}(\Omk)$ is \emph{Kolmogorov-consistent} with
the coarsening system $\Omega_\bullet$ iff
\begin{equation*}
  \Tbstar \, \nu_{k+1} \;=\; \nu_k \qquad \text{for every } k \ge 0.
\end{equation*}
\end{definition}

\begin{theorem}[Consistency of the Wilson family]
\label{thm:wilson-consistency}
The Wilson-measure family $\{\muk\}_{k \ge 0}$ generated by the
UV-direction block-spin flow satisfies the Kolmogorov consistency
condition:
\begin{equation*}
  \Tbstar \, \mu_{k+1} \;=\; \muk \qquad \text{for every } k \ge 0.
\end{equation*}
\end{theorem}

\begin{proof}[Proof sketch]
The block-spin map $\Tb$ is the renormalization-group transformation:
$\Tb$ averages fine-scale link variables over scale-$b$ blocks to
produce coarse-scale link variables, and the renormalized action
$S_k^{\mathrm{Wilson}}$ at scale~$k$ is, by construction, the effective
action obtained from $S_{k+1}^{\mathrm{Wilson}}$ by integrating out
the fluctuations within each block.  This is the defining property of
$\Tb$ as an RG map (cf.\ Bałaban [Bal87]~\S2, Brydges--Federbush--Seiler
[BFS78]).  Hence the pushforward of the fine-scale Gibbs measure under
$\Tb$ is precisely the coarse-scale Gibbs measure:
\begin{equation*}
  \Tbstar \mu_{k+1}(\mathrm{d}\bar U)
  \;=\;
  \int_{\Omkp} \mathbf{1}_{\Tb(U) = \bar U} \, \mu_{k+1}(\mathrm{d}U)
  \;=\;
  \frac{Z_{k+1, \bar U}}{Z_{k+1}} \, \mathrm{d}\bar U
  \;=\;
  Z_k^{-1} \exp\bigl(-S_k^{\mathrm{Wilson}}(\bar U)\bigr) \, \mathrm{d}\bar U
  \;=\;
  \muk(\mathrm{d}\bar U),
\end{equation*}
where $Z_{k+1, \bar U}$ is the partition function at scale~$k+1$
conditioned on the coarse-spin profile $\bar U$.  The block-renormalized
action is defined precisely so that this identity holds; this is the
content of the polymer-expansion convergence proof of
\texttt{formal/sheaf-framework.tex}~\S7 applied at every scale~$k$.
\end{proof}

\begin{remark}[Where the analytic content lives]
\label{rem:where-analysis}
The proof of Theorem~\ref{thm:wilson-consistency} appeals to the
polymer-expansion convergence at every scale, which is in turn
underwritten by the Kotecký--Preiss certificates of the
Phase-2.1~/~3.1--3.3 framework (small-field, large-field, and
buffer-overlap KP).  In the project's current state, those KP
certificates are recorded as
\texttt{SORRY\_LEDGER.md} entries~\texttt{[P2.1]} (analytical content
of \texttt{domGlue\_preserves\_KP}) and the buffer constants of
\texttt{sheaf-framework.tex} Proposition~7.1.  Closing those analytical
sorries closes the proof of Theorem~\ref{thm:wilson-consistency}; no
new analytical content is introduced by Milestone~7.2 \emph{per~se}.
\end{remark}

\section{Tightness via compactness of the target}
\label{sec:tightness}

For the projective-limit theorem to produce a probability measure on
$\Ominf$ (and not merely a finitely-additive cylindrical pre-measure),
the system must be \emph{tight}: for every $\varepsilon > 0$ and every
finite stage~$k$ there must be a compact set $K_\varepsilon \subset
\Omk$ with $\muk(K_\varepsilon) \ge 1 - \varepsilon$.

\begin{proposition}[Trivial tightness]
\label{prop:tightness}
Each $\Omk = SU(2)^{\Ek}$ is itself compact, hence
$K_\varepsilon = \Omk$ works for every $\varepsilon > 0$:
$\muk(\Omk) = 1 \ge 1 - \varepsilon$.  The measure projective system
$\{\muk, \Tbstar\}$ is therefore (trivially) tight.
\end{proposition}

This is the structural payoff of the compactness of $SU(2)$ — the
constructive analog of \emph{boundedness of fields after gauge fixing}
in the noncompact-group case (where tightness becomes the load-bearing
analytical input).  For SU(2), no further analytical work is needed.

\section{The Daniell--Kolmogorov extension theorem and $\muinf$}
\label{sec:extension}

\begin{theorem}[Daniell--Kolmogorov for compact-Hausdorff inverse systems]
\label{thm:dk-extension}
Let $\{X_k, f_{k+1,k}\}_{k \in \mathbb{N}}$ be an inverse system of
nonempty compact Hausdorff spaces and continuous maps, and let
$\{\nu_k\}_{k \in \mathbb{N}}$ be a Kolmogorov-consistent family of
Borel probability measures on the $X_k$.  Then there exists a unique
Borel probability measure $\nu_\infty$ on the projective limit
$X_\infty = \varprojlim X_k$ such that
\begin{equation*}
  (\pi_k)_* \, \nu_\infty \;=\; \nu_k \qquad \forall k.
\end{equation*}
\end{theorem}

This is the classical Daniell--Kolmogorov theorem
(see, e.g., Bourbaki \emph{Intégration} Ch.~IX or Bogachev
\emph{Measure Theory} Vol.~2~\S10.10 for the compact-Hausdorff
formulation; the proof uses the Riesz representation theorem on
$C(X_\infty)$, with the consistent family $\{\nu_k\}$ supplying a
positive linear functional on the dense subalgebra of cylindrical
functions).  Compactness makes both the existence and uniqueness
parts immediate; tightness is automatic
(Proposition~\ref{prop:tightness}).

\begin{definition}[The continuum SU(2) Wilson measure]
\label{def:mu-infinity}
The \emph{continuum SU(2) Wilson measure} $\muinf$ is the unique
Borel probability measure on $\Ominf$ given by
Theorem~\ref{thm:dk-extension} applied to the Wilson-measure
projective system $\{\muk, \Tbstar\}_{k \ge 0}$, which is
Kolmogorov-consistent (Theorem~\ref{thm:wilson-consistency}) and
tight (Proposition~\ref{prop:tightness}).
\end{definition}

\begin{corollary}[Cylindrical realization]
\label{cor:cylindrical}
For every $k$ and every Borel $E \subseteq \Omk$,
$\muinf(\pi_k^{-1}(E)) = \muk(E)$.  In particular, the cylindrical
$\sigma$-algebra $\sigma(\bigcup_k \pi_k^{-1} \mathcal{B}(\Omk))$
generates a sub-$\sigma$-algebra of $\mathcal{B}(\Ominf)$ on which
$\muinf$ is determined by its finite-stage marginals $\{\muk\}$.
\end{corollary}

\section{Injection of $\sighinf > 0$ into $\muinf$}
\label{sec:sigma-injection}

The Phase-7.1 framework theorem
\texttt{sigma\_hat\_uv\_limit\_positive} establishes that the
$Q'$-rational $\sigh_k$ sequence stays bounded below by a strictly
positive $\varepsilon > 0$ at every finite stage:
\begin{equation*}
  \forall N, \quad \sum_{k < N} \eta_k + \varepsilon \;\le\; \sigh_0,
  \qquad \text{i.e.,} \qquad
  \sigh_N \;\ge\; \varepsilon \;>\; 0.
\end{equation*}
The continuum limit $\sighinf$ is well-defined as the (decreasing,
bounded-below) limit of the $\sigh_k$ sequence; it inherits the
strict positivity:
\begin{equation*}
  \sighinf \;\ge\; \varepsilon \;>\; 0.
\end{equation*}

\begin{proposition}[Area-law decay under $\muinf$]
\label{prop:area-law-mu-infinity}
For every rectangular Wilson loop $C \subseteq \Ominf$ with physical
sides $L_1, L_2 \ge a_0$ (so that $C$ is resolved at scale~$0$), the
$\muinf$-expectation of the trace satisfies
\begin{equation*}
  \langle W_R(C) \rangle_{\muinf}
  \;=\;
  \lim_{k \to \infty} \langle W_R(C) \rangle_{\muk}
  \;\le\;
  C_R \, \exp\bigl(-\sighinf \cdot A(C)\bigr),
\end{equation*}
where $A(C) = L_1 L_2$ is the loop area, $\sighinf \ge \varepsilon > 0$
by Phase~7.1, and $C_R$ is the representation-dependent prefactor of
\texttt{docs/outline.md} \S7.1.
\end{proposition}

\begin{proof}
$\langle W_R(C) \rangle_{\muk}$ is bounded by
$C_R \exp(-\sigh_k \cdot A(C))$ at every finite scale~$k$ by the
inductive-step closure of Phase~3 + dimensional scaling
$\sigma_k = a_k^{-2} \sigh_k$.  Taking the limit $k \to \infty$ on
both sides (using Corollary~\ref{cor:cylindrical} to pass the
expectation through to $\muinf$, and using continuity of the
exponential), the bound is preserved with $\sighinf$ in place of
$\sigh_k$.
\end{proof}

\begin{remark}[Why this is the bridge to the mass gap]
\label{rem:bridge-mass-gap}
Phase~8.1 verifies the Osterwalder--Schrader axioms for $\muinf$;
Phase~8.2 applies the OS reconstruction theorem to obtain the
continuum Hilbert space $\mathcal{H}_\infty$ and Hamiltonian
$H_{\mathrm{op}}$.  Phase~9 then identifies the mass gap with the
exponential-decay rate $\sighinf$ via the standard
Wilson-loop$\to$transfer-matrix correspondence: an
area-law-decaying Wilson loop expectation under $\muinf$ implies a
spectral gap $\ge \sighinf$ for $H_{\mathrm{op}}$.
Proposition~\ref{prop:area-law-mu-infinity} is the load-bearing
input that lifts the Q'-rational $\sigh_\infty$ statement into the
continuum measure-theoretic setting where the mass-gap theorem
lives.
\end{remark}

\section{Lean type signature and the paper-only existence decision}
\label{sec:lean-signature}

The Lean module \texttt{lean/YMGap/Continuum/ContinuumMeasure.lean}
provides a type-level interface for the data of this construction:

\begin{itemize}
  \item \texttt{ProjectiveSystem}: a structure capturing the data of
    Definitions~\ref{def:config-system} and~\ref{def:measure-system}
    --- the indexed family of measure spaces and the coarsening
    morphisms, abstract over the (paper-level) consistency and
    tightness conditions.
  \item \texttt{ContinuumMeasure}: a type built from
    \texttt{ProjectiveSystem} plus an abstract token for ``the
    Daniell--Kolmogorov extension exists.''  The extension theorem
    itself is referenced back to Theorem~\ref{thm:dk-extension} of
    this document; no Lean proof is provided.
  \item \texttt{su2\_continuum\_measure}: the SU(2) instance built
    from the Phase-2~/~Phase-3~/~Phase-6 inductive-state machinery.
\end{itemize}

\subsection*{Decision: paper-only existence at Milestone~7.2}

Per the work-plan Milestone~7.2 open question (``decide explicitly
at the kickoff whether to invest in \texttt{Constructive/Reals.lean}
or to accept a Lean-paper hybrid''), this project commits to the
\textbf{Lean-paper hybrid}.  Rationale:

\begin{enumerate}
  \item The Daniell--Kolmogorov extension uses the Riesz representation
    theorem on $C(\Ominf)$, which in a constructive setting requires
    the Bishop-style constructive measure theory of compact Hausdorff
    spaces --- a major analytical build-out far beyond the current
    $Q'$-rational stack.
  \item The mass-gap statement of Phase~9 can be expressed in Lean
    against the type signature of \texttt{ContinuumMeasure}; the
    paper-level existence proof here lets Phase~9 land in Lean without
    waiting on Constructive/Reals.
  \item Milestone~7.3 of the work-plan reserves the option of building
    \texttt{Constructive/Reals.lean} later; this decision does not
    foreclose that.
\end{enumerate}

\section{Open analytical items}
\label{sec:open}

Items deferred from Milestone~7.2 to subsequent milestones:

\begin{itemize}
  \item \textbf{Identification with the OS continuum measure.}
    Section~\ref{sec:sigma-injection} shows $\muinf$ satisfies area-law
    decay; identifying $\muinf$ with the unique OS-axiom-compliant
    continuum SU(2) Wilson measure is Phase~8.1.
  \item \textbf{Reflection positivity for $\muinf$.}  Follows from
    reflection positivity of the finite-stage $\muk$ via the
    projective limit; the lattice argument is standard (lattice link
    variables on each timeslice + reflection through a fixed plane).
    To be carried out as part of Phase~8.1.
  \item \textbf{Euclidean / $O(4)$ invariance.}  At each finite stage
    only the hypercubic symmetry holds; the full $O(4)$ invariance is
    a property of the continuum limit and requires a separate argument
    via the renormalization-group flow restoring rotation invariance.
    Phase~8.1 open item; not addressed here.
  \item \textbf{Exponential clustering.}  Follows from Proposition
    \ref{prop:area-law-mu-infinity} (area-law $\Rightarrow$ exponential
    decay of two-loop correlations for separated loops) via standard
    arguments.  Phase~8.1.
\end{itemize}

\section*{References}

\noindent[Bal87] T.~Bałaban, ``Renormalization group approach to lattice
gauge field theories.~I,'' Comm.\ Math.\ Phys.\ \textbf{109} (1987),
249--301.

\noindent[BFS78] D.~Brydges, J.~Fröhlich, E.~Seiler, ``Some new
results in spectral theory and lattice gauge theories,''
IHES preprint (1978).

\noindent[Bog07] V.~Bogachev, \emph{Measure Theory}, Vol.~2, Springer
2007, \S10.10 (Daniell--Kolmogorov extension for projective systems on
compact Hausdorff spaces).

\noindent[Bou04] N.~Bourbaki, \emph{Intégration} Ch.~IX (projective
limits of measures).

\noindent[JW06] A.~Jaffe, E.~Witten, ``Quantum Yang--Mills theory,''
Clay Millennium Problem statement (2006).
